% BEGIN EMBEDDED BIBLIOGRAPHY
% Bibliography data are part of this source; the compiler writes its temporary input.
\begin{filecontents*}[overwrite]{\jobname.bib}
@unpublished{QED3Master,
 title={The field-strength {Osterwalder--Schrader} axioms for {QED} in three dimensions},
 date={2026-10-20}, note = {Comprehensive edition, 263 pp., 20 October 2026, before the modular integration revision}}
@unpublished{ChamberLattice,
 title={Reflection positivity for chamber {QED} on the infinite lattice},
 date={2026-10-19}, note={Companion manuscript}}
@unpublished{ReferenceResponse,
 title={Exact source identities and normalized reference response for lattice {QED} in three dimensions},
 date={2026-10-20}, note={Companion manuscript}}
@unpublished{TerminalPressure,
 title={The terminal measure and pressure estimates for lattice {QED} in three dimensions},
 date={2026-10-20}, note={Companion manuscript}}
@unpublished{PositiveProbe,
 title={Positive compact probes and free response coefficients for lattice {QED} in three dimensions},
 date={2026-10-20}, note={Companion manuscript}}
@unpublished{ContinuumExistence,
 title={Gaussian domination and structural properties of continuum limits},
 date={2026-10-20}, note={Companion manuscript}}
@unpublished{ContinuumComponents,
 title={Ergodic components and clustering for reflection-positive fields},
 date={2026-10-20}, note={Companion manuscript}}
@unpublished{CommonTuning,
 title={Agreement across covers and common tuning for lattice {QED} in three dimensions},
 date={2026-10-20}, note={Companion manuscript}}
\end{filecontents*}
% END EMBEDDED BIBLIOGRAPHY
% VTM1: volume transfer and component cumulants.
\documentclass[11pt,letterpaper,reqno]{amsart}
% BEGIN INLINE manuscript-style
\makeatletter
\NeedsTeXFormat{LaTeX2e}[2020/10/01]
% Use with \documentclass[11pt,letterpaper,reqno]{amsart}.
% This file controls presentation only. It supplies no mathematical content.

\RequirePackage[T1]{fontenc}
\RequirePackage{lmodern}
\RequirePackage{microtype}
\RequirePackage[letterpaper,margin=1in,headheight=15pt]{geometry}
\RequirePackage{amsmath,amsthm,amssymb,mathtools}
\RequirePackage{enumitem}
\RequirePackage{perpage}
\RequirePackage{xcolor}
\definecolor{manuscriptred}{rgb}{0.6,0.1,0.1}
\definecolor{manuscriptblue}{rgb}{0.1,0.3,0.5}
\definecolor{manuscriptmagenta}{HTML}{60445C}
\RequirePackage[breaklinks=true,colorlinks=true,
  linkcolor=manuscriptred,citecolor=manuscriptblue,urlcolor=manuscriptmagenta]{hyperref}
\RequirePackage{cleveref}
\hypersetup{pdfauthor={}}

% Retain the AMS paragraph, title, heading and running-head conventions.
% Equations retain one global Arabic sequence, as in the reference.
\setcounter{tocdepth}{1}
\renewcommand{\thefigure}{\Roman{figure}}
\renewcommand{\thetable}{\Roman{table}}
\MakePerPage{footnote}
\renewcommand{\thefootnote}{%
  \ifnum\value{footnote}>9\arabic{footnote}\else\fnsymbol{footnote}\fi}

% Upright statements. Main results are A, B, C; supporting statements
% share a section-local sequence. These end marks indicate typography,
% never proof completeness, validation, or a deferred-proof status.
\theoremstyle{definition}
\newtheorem{ms@theorem}{Theorem}
\renewcommand{\thems@theorem}{\Alph{ms@theorem}}
\newtheorem*{ms@theoremstar}{Theorem}
\newtheorem{ms@proposition}{Proposition}[section]
\newtheorem{ms@lemma}[ms@proposition]{Lemma}
\newtheorem{ms@definition}[ms@proposition]{Definition}
\newtheorem{ms@corollary}[ms@proposition]{Corollary}
\newtheorem*{ms@corollarystar}{Corollary}

\newcommand{\ms@majorstart}{%
  \pushQED{\qed}\renewcommand{\qedsymbol}{$\blacksquare\square$}}
\newcommand{\ms@minorstart}{%
  \pushQED{\qed}\renewcommand{\qedsymbol}{$\blacksquare$}}
\newenvironment{theorem}[1][]{\ms@majorstart\ms@theorem[#1]}
  {\popQED\endms@theorem}
\newenvironment{theorem*}[1][]{\ms@majorstart\ms@theoremstar[#1]}
  {\popQED\endms@theoremstar}
\newenvironment{proposition}[1][]{\ms@minorstart\ms@proposition[#1]}
  {\popQED\endms@proposition}
\newenvironment{lemma}[1][]{\ms@minorstart\ms@lemma[#1]}
  {\popQED\endms@lemma}
\newenvironment{definition}[1][]{\ms@minorstart\ms@definition[#1]}
  {\popQED\endms@definition}
\newenvironment{corollary}[1][]{\ms@majorstart\ms@corollary[#1]}
  {\popQED\endms@corollary}
\newenvironment{corollary*}[1][]{\ms@majorstart\ms@corollarystar[#1]}
  {\popQED\endms@corollarystar}
% The source's lemma/proposition variants are deliberately not carried over:
% choose one marker for each statement type rather than encode proof status.

\theoremstyle{remark}
\newtheorem{remark}[ms@proposition]{Remark}
\newtheorem*{remark*}{Remark}
% The ordinary amsthm proof environment retains its open square.

\crefname{ms@theorem}{theorem}{theorems}
\Crefname{ms@theorem}{Theorem}{Theorems}
\crefname{ms@proposition}{proposition}{propositions}
\Crefname{ms@proposition}{Proposition}{Propositions}
\crefname{ms@lemma}{lemma}{lemmas}
\Crefname{ms@lemma}{Lemma}{Lemmas}
\crefname{ms@definition}{definition}{definitions}
\Crefname{ms@definition}{Definition}{Definitions}
\crefname{ms@corollary}{corollary}{corollaries}
\Crefname{ms@corollary}{Corollary}{Corollaries}
\crefname{remark}{remark}{remarks}
\Crefname{remark}{Remark}{Remarks}

% Bibliography configuration for this manuscript.
  \RequirePackage[backend=biber,style=alphabetic,backref=true,
    giveninits=true,isbn=false,date=year]{biblatex}
  \addbibresource{\jobname.bib}
  \setcounter{biburlnumpenalty}{100}
  \setcounter{biburllcpenalty}{100}
  \setcounter{biburlucpenalty}{100}
  \renewcommand*{\bibfont}{\footnotesize}
  \renewbibmacro{in:}{}
  \DeclareFieldFormat{pages}{#1}
  % Use the public alphabetic style. No private author-hash overrides;
  % repeated author names are printed in full.
\makeatother
% END INLINE manuscript-style
\newcommand{\R}{\mathbb R}
\newcommand{\C}{\mathbb C}
\newcommand{\Z}{\mathbb Z}
\newcommand{\E}{\mathbb E}
\newcommand{\dd}{\mathrm d}
\DeclareMathOperator{\supp}{supp}
\DeclareMathOperator{\tr}{tr}
\allowdisplaybreaks
\title[Volume transfer and component cumulants]{Volume transfer and component cumulants for lattice QED in three dimensions}
\date{}
\hypersetup{pdftitle={Volume transfer and component cumulants for lattice QED in three dimensions},pdfauthor={}}
\begin{document}
\begin{abstract}
We transfer normalized reference-volume response estimates for
three-dimensional chamber lattice QED to products of local Wick
polynomials on the prescribed family of arbitrarily large tori.
A two-volume partition comparison controls the normalization.
Its photon contribution loses an explicit, small power of the charge;
Gaussian domination controls the nonzero flux sectors.
We pass the resulting product bounds to any prescribed volume-first
continuum law and then to its translation-ergodic components.
Under the stated reference estimates and component Gaussian bound,
a probe with positive free fourth coefficient has a positive fourth
cumulant in almost every component at sufficiently small charge.
\end{abstract}
\maketitle
\tableofcontents
\section{Introduction}\label{sec:vtm-introduction}

Let $\mathbb P_{n;q,\boldsymbol\mu}^{\rm af}$ be the all-flux
chamber QED law on $C(n)=\Z^3/n\Z^3$, with charge $q>0$,
real lattice mass $\boldsymbol\mu$, Wilson parameter $0<\rho<1$,
weights $(6,6\sqrt2,6\sqrt2,6)$ and antiperiodic fermions.
Its real curvature is $F^{\rm u}=dA+(2\pi/q)z$, where $z$ is
a closed integer two-cochain. Write $Z_n$ for its unnormalized
partition function and $F^\sharp$ for the physical step field:
on the cube of side $\eta$ centered at a plaquette, its
corresponding component is $\eta^{-3/2}F^{\rm u}$.
The measure and its normalization are recalled in
\cref{sec:vtm-data}.

Fix an odd integer $L\ge13$, $M=L^k$ with integer $k\ge2$,
$0<\varepsilon\le1/200$ and $0<\bar m\le1$. Set
\begin{equation}\label{eq:cutoff-scales}
 h=L^{\lfloor\log_L(1/\bar m)\rfloor},\qquad
 \eta=L^{-N},\quad q=e\sqrt\eta,\qquad R=8hL^T,\quad n=R/\eta.
\end{equation}
At each $e,N$ let $\boldsymbol\mu_N$ be real and independent
of the outer volume $T$. Every reference and outer law below
uses this same mass. The cutoffs $N,T$ are nonnegative integers.
For $e\sqrt h<1$, put
\[
 v_K=-\log(e\sqrt h),\qquad
 r_K=\min\{L^i:i\in\Z_{\ge0},\ L^i\ge v_K^2\},
 \qquad R^*=\frac{hMr_K}{L}=hL^{T^*}.
\]
The reference torus has side $n_*=R^*/\eta$ and
$N_*=8L^3$ cells, each of physical width $\ell_c=R^*/(2L)$.

Fix real nonzero $\beta\in C_c^\infty(\R^3;\Lambda^2)$ and
$f=dd^*\beta$. Let $D_\beta$ be the side of a least coordinate
cube containing $\supp\beta$ and let $Q_\beta$, of side
$D_\beta^{\rm buf}=D_\beta+4$ and center $z_\beta$, be that
cube enlarged by two in each direction. Denote cube averaging
by $J_\eta$ and centered differences by
\[
 \delta_{\eta,\mu}g=
 \frac{g(\,\cdot+\eta e_\mu/2)-g(\,\cdot-\eta e_\mu/2)}{\eta}.
\]
Define
\begin{equation}\label{eq:cells-smooth-test}
 a'_{\eta,\nu}=-\sum_{\mu\ne\nu}\delta_{\eta,\mu}\beta_{\mu\nu},
 \qquad
 f^\eta_{\mu\nu}=\delta_{\eta,\mu}a'_{\eta,\nu}
                     -\delta_{\eta,\nu}a'_{\eta,\mu}.
\end{equation}
The covering bond source is
$\widehat a^o_\nu(x)=\eta^{1/2}(J_\eta a'_{\eta,\nu})(\eta x+\eta e_\nu/2)$.
Put $v_\eta=\|d\widehat a^o\|_{\ell^2}^2$.
A centered placement $v\in\eta\Z^3$ satisfies
$|v+z_\beta-(\ell_c/2)(1,1,1)|_\infty\le\eta/2$.
Translate the probe by $v$, reflect its source into the
$N_*$ cells, and periodize to obtain $\widehat a_c$.
The reflected copies use the two-form orientation signs.

Let $D^{\rm ch}[B]$ be the chamber Dirac operator with bond
phase $B$. Define the free one-cell ratio and its ordinary
derivatives by
\begin{equation}\label{eq:response-coefficient-relations}
 N_0(t)=\frac{\det(D^{\rm ch}[q(t/e)\widehat a_0]+\bar m\eta)}
                    {\det(D^{\rm ch}[0]+\bar m\eta)},\qquad
 c_m=N_0^{(m)}(0),\quad m=2,4.
\end{equation}
The free mass is positive even when $\boldsymbol\mu_N$ is signed.
Write $\mathbb E^{(0)}_{n_*}$ for the same-mass reference law
conditioned on zero flux, and set
\begin{equation}\label{eq:vtm-Wick}
 H_2(x;v)=x^2-v,\qquad H_4(x;v)=x^4-6vx^2+3v^2,\qquad
 W_m(x;v,c)=e^{-m}H_m(x;v)-c.
\end{equation}
Here and below $e$ is the charge, so $e^{-m}$ is its inverse
power. With $X_c=\langle d\widehat a_c,F^{\rm u}\rangle$, put
\begin{equation}\label{eq:vtm-reference-root}
 \Theta_m^{(0)}
 =\left|\mathbb E_{n_*}^{(0)}
       \prod_{c\in\{0,\ldots,2L-1\}^3}W_m(X_c;v_\eta,c_m)
   \right|^{1/N_*}.
\end{equation}
Let $z_{r,n}$ be the positive zero-flux matter factor in
\eqref{eq:external-matter-factor}. At continuum mass $\bar m$ define
\[
 b(f)=-6\tr\left[
    \bigl((\bar m+\gamma\cdot\partial)^{-1}i\gamma\cdot d^*\beta
                 \bigr)^4\right],
\]
where $\gamma_\nu=\operatorname{diag}(\sigma_\nu,-\sigma_\nu)$.
The displayed fourth power is trace class and its trace is real
by \cite[Theorem~A]{PositiveProbe}.

\begin{theorem}[Volume transfer and component cumulants]\label{thm:vtm-main}
Fix the data above and a charge interval $0<e\le e_*\le1/3$.
For each charge let $N_{\rm ref}(e)$ be a finite nonnegative integer.
Assume the following conditions for every charge in this interval
and every $N\ge N_{\rm ref}(e)$.
\begin{enumerate}[label=\textup{(\roman*)}]
\item The scale and placement conditions are
\[
 v_K\ge8,\qquad e\sqrt h\log h\le1,\qquad
 6N_*Mv_K^2\exp(-v_K)
          \bigl(v_K\exp(-v_K)+\tfrac12\bigr)\le\pi^2,
\]
and
\begin{equation}\label{eq:cells-fit}
 D_\beta^{\rm buf}+3\eta\le\ell_c,\qquad
 v+Q_\beta\subset[\ell_c/8,7\ell_c/8]^3
\end{equation}
for each centered placement used.
\item There are finite nonnegative constants $C_{\rm mat}$ and
$C_{\rm ref}$, independent of charge, mesh and placement, such that,
at the same $q,\boldsymbol\mu_N$ on the two reference tori,
\begin{align}
 L^3\log z_{r,n_*}-\log z_{r,Ln_*}&\le C_{\rm mat},\notag\\
 \Theta_4^{(0)}&\le C_{\rm ref}e^{\mathfrak e},\qquad
 (1+|c_2|)\Theta_2^{(0)}\le C_{\rm ref}e^{\mathfrak e},
 \label{eq:vtm-reference-hypotheses}\\
 |c_2|+|c_4|^{1/2}&\le C_{\rm ref},\qquad
 |c_4-3c_2^2-b(f)|\le C_{\rm ref}(e+\eta^{1/4}),\notag
\end{align}
where
\[
 \mathfrak e=\tfrac12(1/4-9\varepsilon)-12\varepsilon,\qquad
 a_L=\frac{L^3-1}{(L-1)^3},\qquad
 \sigma=\mathfrak e-\frac{2a_L}{N_*}>0.
\]
\end{enumerate}
There is $C_U<\infty$, depending only on these fixed data and
constants, such that the following assertions hold uniformly
in the stated charge interval.

For each $N\ge N_{\rm ref}(e)$ there are $u_2,u_4\ge0$,
independent of the outer cutoff $T$, with
\begin{equation}\label{eq:vtm-main-errors}
 u_4\le C_Ue^\sigma,\qquad
 (1+|c_2|)u_2\le C_Ue^\sigma,\qquad
 |c_4-3c_2^2-b(f)|\le C_U(e^\sigma+\eta^{1/4}).
\end{equation}
For every $T\ge T^*(e)+1$ and every finite
$Y\subset(R^*/L)\Z^3$ whose points are incongruent modulo $R\Z^3$,
\begin{equation}\label{eq:vtm-main-products}
 \left|\E_{N,T}\prod_{y\in Y}
 W_m(F^\sharp(\tau_yf^\eta);v_\eta,c_m)\right|
              \le u_m^{|Y|},\qquad m=2,4.
\end{equation}
Here $\tau_yf(x)=f(x-y)$ and $\E_{N,T}$ is expectation under
$\mathbb P_{n;q,\boldsymbol\mu_N}^{\rm af}$. Its smooth-test
fourth cumulant satisfies
\begin{equation}\label{eq:vtm-main-finite-cumulant}
 |\kappa_4(F^\sharp(f))-e^4b(f)|
 \le(8C_U+3C_U^2)e^{4+\sigma}
                  +C_U(\eta^2+e^4\eta^{1/4}).
\end{equation}

Fix one charge. Let $\mu$ be any prescribed volume-first limit:
for $N_i\to\infty$, choose fixed-spacing weak volume limits
$\mu_i$ of these laws, with $\mu_i\Rightarrow\mu$.
All weak convergence here is on the negative Hermite space
$H_{-4}$ specified in \cref{sec:continuum-limits}.
Assume that $\mu$ is translation invariant and that almost every
component in its translation-ergodic decomposition
$\mu=\int\mu_\xi\,\dd\nu(\xi)$ satisfies
\begin{equation}\label{eq:vtm-component-G}
 \E_{\mu_\xi}\exp(F(g))\le\exp(\|g\|_2^2/2)
 \quad\text{for all real }g\in\mathcal S(\R^3;\Lambda^2).
\end{equation}
Then, on one $\nu$-full set,
\begin{equation}\label{eq:vtm-main-component}
 |\kappa_4^{\mu_\xi}(f)-e^4b(f)|
                    \le C_fe^{4+\sigma},\qquad C_f=8C_U+3C_U^2.
\end{equation}
If $b(f)>0$ and $C_fe^\sigma\le b(f)/2$, every component
in this set has $\kappa_4^{\mu_\xi}(f)\ge e^4b(f)/2>0$.
The full-measure set is for this fixed charge, probe and parent law.
\end{theorem}

For the tuned chamber construction, the same mass assignment
is selected on the top torus by \cite[Theorem~A]{CommonTuning}.
The reference-response estimate is \cite[Theorem~A]{ReferenceResponse},
and the two-torus matter comparison is \cite[Theorem~A]{TerminalPressure}.
The free coefficients and a probe with $b(f)>0$ are supplied by
\cite[Theorems~A--B]{PositiveProbe}. The volume-first laws are
constructed in \cite{ContinuumExistence}; the component Gaussian
bound is the last clause of \cite[Theorem~A]{ContinuumComponents}.
The simultaneous admission of the construction parameters is
separate from the implication proved here.

The proof is the volume-transfer and component argument of
\cite[\S25--26]{QED3Master}, with its finite-periodic moment
passage from \cite[\S6.1--6.2]{QED3Master}.
Its first difficulty is normalization. Folding bounds the
unnormalized observable by its value on the small reference torus,
but the denominator belongs to the larger torus. Sidewise
convexity of the partition function controls this denominator
using the two reference volumes. The exact photon determinant
then shows that the normalization costs only the power
$e^{-2a_L/N_*}$; the response estimate retains the positive
exponent $\sigma$.

The reference response is in flux zero. We split the full
reference expectation by flux and bound its nonzero-flux part
with the Gaussian tail and high moments. This gives bounds on
every finite product of distinct translated Wick polynomials,
uniformly in the outer volume. Their passage to a prescribed
parent uses the periodic seminorm at finite volume and the sharp
$L^2$ norm only after the supports fit in a period cube.
Finally, spatial averages and a count of repeated indices bound
the invariant conditional expectations of both Wick polynomials.
These two conditional bounds determine the fourth cumulant in
each component.

The charge restriction for positivity is imposed after $C_U$:
\begin{equation}\label{eq:constants-probe-charge}
 0<e_\beta\le
       \min\{e_*,[b(f)/(2C_f)]^{1/\sigma}\}.
\end{equation}
It is positive when $b(f)>0$ and $C_f>0$.
The conclusions concern subsequential limits of the stated
cutoff family; no uniqueness of the continuum law is required.

\section{The chamber law and the reference data}\label{sec:vtm-data}

\subsection{Normalization, flux and Gaussian domination}

The chamber operator is the finite-range operator of
\cite[\S2]{ChamberLattice}; its bond argument is a phase.
It acts on four spin components and $48$ chamber labels per
site. The fermions are antiperiodic in each coordinate.
For a rectangular coordinate torus with $V$ sites and every
side at least four, put $r_q=2\pi/q$. The unnormalized integral
of a curvature observable $O$ is
\begin{equation}\label{eq:vtm-law}
 Z[O]=\sum_{\substack{z\in C^2(\Z)\\dz=0}}
  \int_{[0,r_q)^{3V}}
 O(dA+r_qz)\exp\bigl(-\tfrac12\|dA+r_qz\|^2\bigr)
 \det(D^{\rm ch}[qA]+\boldsymbol\mu)\,\dd A.
\end{equation}
Cochains and the integer sum are on this torus; norms and
pairings count positively oriented cells. The identification
$(A,z)\sim(A+r_qk,z-dk)$ makes the integral independent
of the chosen fundamental domain. The law is $Z[O]/Z[1]$.
The partition function is finite and positive for every real
$\boldsymbol\mu$, and the determinant gives nonnegative weights;
see \cite[\S2--3]{ChamberLattice}. Flux is the three integer
periods of $z$, ordered as $(23,13,12)$.
We retain every flux sector in \eqref{eq:vtm-law}.

The coordinate-torus Gaussian bound
\cite[Proposition~5.1]{ChamberLattice} is
\begin{equation}\label{eq:slab-gaussian-domination}
 \E^{\rm af}\exp\langle g,F^{\rm u}\rangle
       \le\exp(\|g\|^2/2)
\end{equation}
for every real torus two-cochain $g$.
It holds for both odd and even coordinate sides.
The law is invariant under lattice translations and
$F^{\rm u}\mapsto-F^{\rm u}$. In particular all odd
moments of a real linear observable vanish.

On $C(n)$, let $\gamma_n$ be the standard Gaussian measure
on the exact real two-cochains, and let $A(u)$ be the
mean-zero coclosed potential of $u$.
For $\theta\in[-\pi,\pi]^3$ put
$H_{\theta,\nu}=\theta_\nu/(qn)$. The matter factor is
\begin{equation}\label{eq:external-matter-factor}
 z_{r,n}=\int_{[-\pi,\pi]^3}\frac{\dd\theta}{(2\pi)^3}
     \int\det(D^{\rm ch}[q(A(u)+H_\theta)]
                     +\boldsymbol\mu)\,\dd\gamma_n(u)>0.
\end{equation}
The subscript $r$ denotes the regulator. The exact zero-flux
factorization of \cite[\S3.4]{ReferenceResponse} gives
\begin{align}
 Z_n^{(0)}[1]&=Z_n^{\rm ph}z_{r,n},
                     \label{eq:external-zero-factorization}\\
 Z_n^{\rm ph}&=r_q^{\,n^3+2}(2\pi)^{n^3-1}
                         \exp(-\mathcal L(n)/2).
                     \label{eq:external-photon-factor}
\end{align}
where $\mathcal L(n)$ is the scalar-Laplacian logarithmic
determinant in \cref{sec:vtm-defect}.
Thus the normalized zero-flux curvature law is the Gaussian
law weighted by the twist-averaged determinant in
\eqref{eq:external-matter-factor}. The factorization retains
the gauge and harmonic volume factors; they will contribute
to the partition defect.

\subsection{The source and cell contract}\label{sec:reference-cells}

The source identities and estimates in this subsection are
proved in \cite[\S2]{ReferenceResponse}. We state the needed
contract with its support and normalization.

Write $p_{\rm cell}=n_*/L$ and $w=p_{\rm cell}/2$.
The reference side $n_*$ and $p_{\rm cell}$ are odd,
so $w$ is half-integral. The number $2L$ of cells per
side is even. In unit-lattice coordinates the open cells
are $Q_c^\circ=w(c+(0,1)^3)$.

For a two-form $g$ let
\[
 g_{[\eta]}(\sigma)=\eta^{-3/2}
        \int_{Q_\eta(\sigma)}g_{\mu\nu}(y)\,\dd y,\qquad
 Q_\eta(\sigma)=\eta\bigl(x+(e_\mu+e_\nu)/2+[-1/2,1/2)^3\bigr),
\]
when $\sigma=(x;\mu\nu)$. The source of the introduction obeys
$d\widehat a^o=(f^\eta)_{[\eta]}$.
After periodizing, this gives
\[
 \langle d\widehat a,F^{\rm u}\rangle=F^\sharp(f^\eta).
\]
All $f^\eta$ are supported in $Q_\beta$ for $\eta\le1$, and
\begin{equation}\label{eq:cells-test-error}
 \|f^\eta-f\|_{H^s}\le\eta^2\|\beta\|_{H^{s+4}}
                       \quad(s\in\R).
\end{equation}
The variance converges with the quantitative estimate
\begin{equation}\label{eq:cells-variance-limit}
 \left|\sqrt{v_\eta}-\|f\|_2\right|
 \le21\eta|\beta|_{C^3}(D_\beta^{\rm buf})^{3/2}
                       +\eta^2\|\beta\|_{H^4}.
\end{equation}
Here $|\beta|_{C^3}$ is the maximum component derivative norm
through order three. These estimates follow from the centered
difference symbols and the cell-average approximation.

For the placed probe, the coordinate map $g_c$ carrying the
base cell to cell $c$ is, in coordinate $\mu$,
$x_\mu\mapsto x_\mu+2iw$ if $c_\mu=2i$, and
$x_\mu\mapsto(2i+2)w-x_\mu$ if $c_\mu=2i+1$.
Transport oriented cochains by this map before periodizing.
Denote the placed and reflected covering source by
$\widehat a_c^o$; its periodization is $\widehat a_c$.
On every carrying torus with $p_{\rm cell}\mid n$ and $n\ge n_*$,
\begin{equation}\label{eq:cells-orthogonality}
 \langle d\widehat a_c,d\widehat a_{c'}\rangle_n
                         =v_\eta\mathbf1_{\{c=c'\}}.
\end{equation}
The placement condition \eqref{eq:cells-fit} ensures that their
support boxes are strictly inside the open cells and that the
boxes enlarged by one embed in the torus.
Each polynomial $P(X_c)$ is consequently read by the same
local scalar integrand on every such torus.

Here is the scalar observable class used in the transfer,
with the conventions of \cite[\S4.1, (46)--(50)]{ChamberLattice}.
For a coordinate box $B$, let $B^{(j)}$ be its contained
oriented lattice $j$-cells. A scalar local integrand is a Borel
gauge-invariant function $G$ of the real bonds in $B^{(1)}$
such that, for some $C<\infty$ and $\epsilon_B>0$,
\[
 |G(\widehat A)|\le C\exp\left[
   (\tfrac12-\epsilon_B)
          \sum_{\sigma\in B^{(2)}}|d\widehat A(\sigma)|^2\right].
\]
Gauge invariance means $G(\widehat A-d\lambda)=G(\widehat A)$.
A local lift of $(A,z)$ is a real bond cochain $\widehat A$
on $B$ with $d\widehat A=F^{\rm u}$ and
$\exp(iq\widehat A(b))=\exp(iqA(b))$.
The cited construction proves existence and independence of
the lift for these evaluations.

For an open set $U$, let $\mathcal A(U)$ be the linear span
of products $\Phi(d\widehat A)\prod_jG_j(\widehat A|_{B_j})$,
where the boxes $B_j\subset U$ are pairwise disjoint and
$\Phi$ is a bounded Borel function of finitely many curvature
values in $U$. On a quotient with periods $\Gamma$ the boxes
must embed: each family $B_j+[-1,1]^3+\gamma$ is disjoint,
and all $B_j+\gamma$ for distinct $(j,\gamma)$ are disjoint.

For the present probe take a supporting coordinate box
$B_0\Subset Q_0^\circ$ as in \cite[\S2.2]{ReferenceResponse}
and $B_c=g_cB_0$. The local expression for $P(X_c)$ is
$P(\sum_{\sigma\subset B_c}d\widehat a_c^o(\sigma)
                                      d\widehat A(\sigma))$.
Cauchy--Schwarz bounds its argument by
$\sqrt{v_\eta}(\sum_{\sigma\subset B_c}|d\widehat A(\sigma)|^2)^{1/2}$,
so it has a bound
$C_P\exp(\frac14\sum_{\sigma\subset B_c}|d\widehat A(\sigma)|^2)$.
Thus it belongs to $\mathcal A(Q_c^\circ)$.
The reflection is antilinear on coefficients; it fixes the
coefficients of these real Wick polynomials.

Even cell indices give the precise translation formula
\begin{equation}\label{eq:cells-translation-observable}
 c(y)=\frac{2y}{p_{\rm cell}\eta}\in2\Z^3,\qquad
 X_{c(y)}=F^\sharp(\tau_{v+y}f^\eta),
 \quad y\in(R^*/L)\Z^3.
\end{equation}
Two such cells coincide on the torus exactly when the
translations agree modulo its physical periods. Lattice
translation invariance removes the common placement $v$
from their product expectation.

\subsection{The normalized source identity}

On the reference torus define $a(t)=e^{-1}\sum_ct_c\widehat a_c$ and
\begin{equation}\label{eq:external-reference-partition}
 \mathring\Phi(t)=z_{r,n_*}^{-1}
 \int\frac{\dd\theta}{(2\pi)^3}\int
 \det(D^{\rm ch}[q(A(u)+H_\theta+a(t))]
                         +\boldsymbol\mu_N)\,\dd\gamma_{n_*}(u).
\end{equation}
The $\theta$ integral is over $[-\pi,\pi]^3$.
The source-shift identity of
\cite[\S3.6, Proposition~3.2 and (53)]{ReferenceResponse},
valid for arbitrary real lattice mass, is
\begin{equation}\label{eq:external-Wick-product}
 \E^{(0)}_{n_*}\prod_c W_m(X_c;v_\eta,c_m)
       =\left[\prod_c(\partial_{t_c}^m-c_m)\right]\mathring\Phi(0),
       \qquad m=2,4.
\end{equation}
Thus \eqref{eq:vtm-reference-root} is exactly the response norm
controlled there. Its reference law is normalized, and its
coefficients are the ordinary free derivatives
\eqref{eq:response-coefficient-relations}. The vacuum parameter
of a common tuned trajectory cancels in this ratio. In a
two-volume logarithmic defect its extensive vacuum contributions
cancel as well.

\section{Transfer from the reference volumes}\label{sec:uniform-transfer}

The renormalization estimates are obtained on two reference tori.
To use them for volume-first limits, we need bounds on products
of local observables in arbitrarily large cutoff tori. The
chessboard estimate reduces each such product to a folded
one-cell observable. A separate convexity estimate controls
the normalization by the partition functions of the two
reference volumes.

\subsection{Folded observables and the partition defect}

Use unit-lattice lengths in this subsection. Fix the charge $q>0$
and a real lattice mass. Choose coordinate cells of width
$w\in\tfrac12\Z_{>0}$ and offsets $o_\mu\in\tfrac12\Z$:
\[
 Q_c^\circ=\{y:o_\mu+wc_\mu<y_\mu<o_\mu+w(c_\mu+1),\
                           \mu=1,2,3\}.
\]
The walls are admissible coordinate mirrors. For even positive
integers $s_\mu$ with $ws_\mu\ge4$, let $C(ws)$ be the rectangular
torus with these side lengths and antiperiodic fermions.
It has $s_1s_2s_3$ cells. The wall-reflection group acts simply
transitively on the covering cells. Write $O^{[c]}$ for the
reflected copy of $O\in\mathcal A(Q_0^\circ)$ in $Q_c^\circ$,
and
\[
 O^\sharp_{ws}=\prod_{c\bmod s}O^{[c]},\qquad
 \mathcal Z_s(O)=Z_{C(ws)}[O^\sharp_{ws}].
\]
The copies use the antilinear reflection action, including
conjugation of scalar coefficients.

These arrays satisfy the rectangular chessboard estimate
\cite[(69)]{ChamberLattice}. The cell theorem is stated there for integer width.
For the coordinate cells used here, its fold proof also applies
at half-integer width: all walls remain admissible, paired
reflections translate by the integer $2w$, and the even cell
counts give integer torus sides at least four. For each layer, the fold
identity is applied with the trace based at that layer's own lower boundary.
For a support box $B=y+[-r,r]^3\subset Q_c^\circ$ one has
$2r<w$. A nonzero period has sup norm at least $2w>4r$
and at least four. Thus its enlarged boxes $B+[-1,1]^3$
are disjoint, by the same embedding estimate used in the
proof of \cite[Proposition~4.3]{ChamberLattice}.
The products therefore belong to the required slab classes.

Hold two entries of $s$ fixed and vary $s_\mu$.
The product of the copies in the layer $c_\mu=0$ gives
a Hilbert--Schmidt slab operator $Y$ independent of $s_\mu$.
Its reflected layer pairs fill the entire folded array.
The slab-fold identity \cite[Proposition~4.2, (66)]{ChamberLattice} gives
\begin{equation}\label{eq:transfer-fold-trace}
 \mathcal Z_s(O)=\operatorname{Tr}(Y^*Y)^{s_\mu/2}.
\end{equation}
In particular $\mathcal Z_s(O)\ge0$, and
$s_\mu\mapsto\mathcal Z_s(O)^{1/s_\mu}$ is nonincreasing.
If $Y\ne0$, the logarithm of this folded partition function
is convex in $s_\mu$. The first assertion is the monotonicity
of Schatten norms, and the second follows from H\"older's
inequality for the trace of positive powers.
The normalization of $Y$ here is that of the slab operators
in \cite[\S4.3]{ChamberLattice}.

Choose an even reference cell count $s_0\ge2$ with
$n_r=ws_0\ge4$, and an integer $j\ge2$. Put
\begin{equation}\label{eq:transfer-reference-data}
 N_r=s_0^3,\qquad
 D_j=j^3\log Z_{C(n_r)}[1]-\log Z_{C(jn_r)}[1],
 \qquad \delta_j=\frac{D_j}{(j-1)^3N_r}.
\end{equation}
For an observable in the base cell its normalized reference
seminorm is
\begin{equation}\label{eq:transfer-reference-seminorm}
 \|O\|_*=
 \left(\frac{\mathcal Z_{(s_0,s_0,s_0)}(O)}
                    {Z_{C(n_r)}[1]}\right)^{1/N_r}.
\end{equation}
The partition function is positive, and the numerator is
nonnegative, so this is well-defined also when it vanishes.

\begin{proposition}\label{prop:transfer-general}
One has $D_j\ge0$. Let $s$ be even, with either $s=s_0$
or $s\ge js_0$, and put $r=s/s_0$, $n=ws$.
For every set $\mathcal C$ of distinct cells of $C(n)$ and
$O_c\in\mathcal A(Q_0^\circ)$,
\begin{equation}\label{eq:transfer-general}
 \left|\E_{C(n)}\prod_{c\in\mathcal C}O_c^{[c]}\right|
                    \le\prod_{c\in\mathcal C}\exp(\delta_j)\|O_c\|_*.
\end{equation}
The right side uses only the two reference partition functions
and the one-cell observables folded on $C(n_r)$.
\end{proposition}

\begin{proof}
Successively decreasing the three side counts in
\eqref{eq:transfer-fold-trace} gives, for $s\ge s_0$,
\begin{equation}\label{eq:transfer-fold-upper}
 \mathcal Z_{(s,s,s)}(O)
            \le\mathcal Z_{(s_0,s_0,s_0)}(O)^{r^3}.
\end{equation}
Taking $O=1$ and $r=j$ proves $D_j\ge0$.

We next obtain the lower bound for the normalizing partition
function. Define
$p(r_1,r_2,r_3)=\log Z_{C(r_1n_r,r_2n_r,r_3n_r)}[1]$
on ratios with integer side lengths at least four.
The pressure trace \cite[(75)]{ChamberLattice} shows that
$p$ is convex in each argument and that $p/r_\mu$ is
nonincreasing in that argument. For $r\ge j$ put
$a=(r-1)/(j-1)$. Convexity gives
$p(r,\cdot,\cdot)\ge a p(j,\cdot,\cdot)
-(a-1)p(1,\cdot,\cdot)$.
Apply this in each argument, retaining the positive coefficient
$a$ at the next step. Bound the three subtracted terms by the
side monotonicity. The result is
\begin{align*}
 p(r,r,r)\ge{}&
 a^3p(j,j,j)\\
 &-(a-1)(a^2j^2+ajr+r^2)p(1,1,1)\\
 ={}&r^3p(1,1,1)-a^3D_j.
\end{align*}
The last equality uses $aj-(a-1)=r$ three times.
The negative coefficients are why upper bounds are used on
the three subtracted terms; no sign condition on $p$ is needed.
Thus, including the equality case $r=1$,
\begin{equation}\label{eq:transfer-partition-lower}
 \log Z_{C(n)}[1]\ge
 r^3\left(\log Z_{C(n_r)}[1]-\frac{D_j}{(j-1)^3}\right).
\end{equation}

Fill the cells outside $\mathcal C$ with the constant integrand.
The normalized rectangular chessboard estimate then gives
\[
 \left|\E_{C(n)}\prod_{c\in\mathcal C}O_c^{[c]}\right|
 \le\prod_{c\in\mathcal C}
       \left(\frac{\mathcal Z_{(s,s,s)}(O_c)}
                         {Z_{C(n)}[1]}\right)^{1/(r^3N_r)}.
\]
Use \eqref{eq:transfer-fold-upper} in the numerator and
\eqref{eq:transfer-partition-lower} in the denominator.
Each factor is at most $\exp(\delta_j)\|O_c\|_*$, proving
\eqref{eq:transfer-general}.
\end{proof}

\section{The photon, matter and flux defect}\label{sec:vtm-defect}

\subsection{The scalar-Laplacian determinant}

The zero-flux Gaussian factor involves the determinant of the
unit-lattice scalar Laplacian with its constant mode removed.
Write its logarithm as
\[
 \mathcal L(n_1,n_2,n_3)
  =\sum_{\substack{k\in\prod_\mu\Z/n_\mu\Z\\k\ne0}}
      \log\left(\sum_{\mu=1}^3
                       \left(2-2\cos\frac{2\pi k_\mu}{n_\mu}\right)\right),
 \qquad \mathcal L(n)=\mathcal L(n,n,n).
\]
For $j\ge2$ put
\[
 G_j=\left(1+\frac j{3.2}\right)\frac{2j^2}{2.56},\qquad
 K_j=(1+j+j^2)(\log j+jG_j).
\]
The estimate used for the two reference volumes is
\begin{equation}\label{eq:transfer-photon-logdet}
 \frac12\bigl(\mathcal L(jn)-j^3\mathcal L(n)\bigr)
              \le-(j^3-1)\log n+K_j
                   \qquad(n,j\in\Z,\ n,j\ge2).
\end{equation}
Its logarithmic term cancels the corresponding scale dependence
of the Gaussian normalization. We retain the finite-volume
calculation since a bulk determinant bound would lose this term.

Sum first in the first coordinate. For a nonzero transverse
momentum $k'=(k_2,k_3)$ define $\omega(k')>0$ by
\[
 2\cosh\omega(k')
     =2+\sum_{\mu=2,3}
                 \left(2-2\cos\frac{2\pi k_\mu}{n_\mu}\right).
\]
Factoring over the $n_1$-th roots of unity gives
\[
 \prod_{k_1\in\Z/n_1\Z}
       \left(2\cosh\omega-2\cos\frac{2\pi k_1}{n_1}\right)
                       =4\sinh^2(n_1\omega/2).
\]
For $k'=0$, omit $k_1=0$; the corresponding product is $n_1^2$.
Therefore
\begin{equation}\label{eq:transfer-photon-one-direction}
 \mathcal L(n_1,n_2,n_3)
   =2\log n_1+
       2\sum_{k'\ne0}\log\bigl(2\sinh(n_1\omega(k')/2)\bigr).
\end{equation}
For $z>0$ let $g(z)=-\log(1-\exp(-2z))$.
The identity $2\sinh z=\exp(z)(1-\exp(-2z))$ shows that
\[
 \log(2\sinh(jz))-j\log(2\sinh z)\le jg(z).
\]
Consequently
\begin{multline}\label{eq:transfer-photon-side-difference}
 \mathcal L(jn_1,n_2,n_3)-j\mathcal L(n_1,n_2,n_3)\\
 \le2\log(jn_1)-2j\log n_1
                       +2j\sum_{k'\ne0}g(n_1\omega(k')/2).
\end{multline}

In the three applications below, $n_1=n$ and
$n_2,n_3\in\{n,jn\}$. Represent each transverse momentum by
$|k_\mu|\le n_\mu/2$ and put $p_\mu=2\pi k_\mu/n_\mu$.
For
$M(p)=\sum_{\mu=2,3}(2-2\cos p_\mu)$,
\[
 \frac4{\pi^2}|p|^2\le M(p)\le8,\qquad
 \omega=2\operatorname{arcsinh}(\sqrt M/2)\ge0.51|p|.
\]
The last inequality follows from concavity of
$\operatorname{arcsinh}$ on $[0,\sqrt2]$ and
$\operatorname{arcsinh}\sqrt2>0.81\sqrt2$.
Thus $n\omega(k')/2\ge c_j|k'|$, where $c_j=1.6/j$.
Since
$g(z)\le\exp(-2z)(1+(2z)^{-1})$, and
$|k'|\ge(|k_2|+|k_3|)/\sqrt2$, one obtains
\begin{align*}
 \sum_{k'\ne0}g(n\omega(k')/2)
 &\le\left(1+\frac1{2c_j}\right)
   \left[
     \left(\sum_{i\in\Z}\exp(-\sqrt2c_j|i|)\right)^2-1
   \right]\\
 &=\left(1+\frac1{2c_j}\right)
                  \sinh(c_j/\sqrt2)^{-2}
 \le\left(1+\frac1{2c_j}\right)\frac2{c_j^2}
 =G_j.
\end{align*}
Enlargement to all of $\Z^2\setminus\{0\}$ is valid because
the chosen representatives are distinct.

Finally telescope the enlargement of the three sides:
\begin{align*}
 \mathcal L(jn,jn,jn)-j^3\mathcal L(n,n,n)
 ={}&\bigl[\mathcal L(jn,jn,jn)-j\mathcal L(n,jn,jn)\bigr]\\
 &+j\bigl[\mathcal L(jn,n,jn)-j\mathcal L(n,n,jn)\bigr]\\
 &+j^2\bigl[\mathcal L(jn,n,n)-j\mathcal L(n,n,n)\bigr].
\end{align*}
Symmetry of $\mathcal L$ supplies this identity.
Each bracket is bounded by
$2\log j+2(1-j)\log n+2jG_j$ by
\eqref{eq:transfer-photon-side-difference}.
Multiplication by $1,j,j^2$ and summation prove
\eqref{eq:transfer-photon-logdet}.

\subsection{The photon, matter and flux contributions}

Fix $q>0$ and a real lattice mass. On a cubic torus with
$n\ge4$ sites per side, let $Z_n^{(w)}[1]$ be the
partition function in flux sector $w\in\Z^3$, and put
$P_n(w)=Z_n^{(w)}[1]/Z_n[1]$. The flux order is
$(23,13,12)$. Define
\begin{equation}\label{eq:transfer-flux-tail}
 \varsigma_n=\sum_{w\ne0}
              \exp\left(-\frac{2\pi^2|w|_2^2}{q^2n}\right).
\end{equation}
The notation $z_{r,n}$ below is the positive matter factor
of \eqref{eq:external-matter-factor}.

\begin{proposition}\label{prop:transfer-flux-defect}
For $n\ge4$ and an integer $j\ge2$,
\begin{align}
 Z_n[1]&=\frac{Z_n^{\mathrm{ph}}z_{r,n}}{P_n(0)},
                  \qquad 0<P_n(0)\le1,\label{eq:transfer-zero-factor}\\
 P_n(w)&\le\exp\left(-\frac{2\pi^2|w|_2^2}{q^2n}\right),
 \qquad P_n(0)\ge1-\varsigma_n. \label{eq:transfer-flux-weight}
\end{align}
If $q^2n\le1$, then
\begin{equation}\label{eq:transfer-flux-small}
 \varsigma_n\le7\exp\left(-\frac{2\pi^2}{q^2n}\right)<\frac18.
\end{equation}
The defect $D_j$ of \eqref{eq:transfer-reference-data}
decomposes as
\begin{align}
 D_j&=D_j^{\mathrm{ph}}+D_j^{\mathrm{mat}}+D_j^{\mathrm{fl}},
                                      \label{eq:transfer-defect-split}\\
 D_j^{\mathrm{mat}}&=j^3\log z_{r,n}-\log z_{r,jn},
 \qquad
 D_j^{\mathrm{fl}}=-j^3\log P_n(0)+\log P_{jn}(0),\notag\\
 D_j^{\mathrm{ph}}
   &=(j^3-1)\log\frac{2\pi}{q^2}
                +\tfrac12[\mathcal L(jn)-j^3\mathcal L(n)]\notag\\
   &\le(j^3-1)\log\frac{2\pi}{q^2n}+K_j.
                                      \label{eq:transfer-photon-defect}
\end{align}
If $\varsigma_n<1$, then
\begin{equation}\label{eq:transfer-flux-defect}
 D_j^{\mathrm{fl}}\le-j^3\log(1-\varsigma_n).
\end{equation}
All identities use the same charge and lattice mass at
the two sides; the mass can have either sign.
\end{proposition}

\begin{proof}
The zero-flux factorization \eqref{eq:external-zero-factorization}
gives $Z_n^{(0)}[1]=Z_n^{\mathrm{ph}}z_{r,n}>0$.
All sectors have nonnegative weights, and their absolutely
convergent sum is $Z_n[1]$. Hence
$0<P_n(0)\le1$ and \eqref{eq:transfer-zero-factor}.

For a configuration of flux $w$, let $h_w$ be the
constant two-cochain with value
$r_qw_{\mu\nu}/n^2$ on plaquettes of orientation $\mu\nu$.
The Fourier Hodge decomposition writes
$r_qz=dB+h$ with $h$ constant.
Summing over a coordinate two-torus annihilates $dB$
and gives $n^2h_{\mu\nu}=r_qw_{\mu\nu}$.
Thus $F-h_w$ is exact. A constant cochain is orthogonal
to every exact cochain, by periodic summation by parts.
In particular $\langle h,F\rangle=\langle h,h_w\rangle$
on sector $w$ whenever $h$ is constant.

Apply Gaussian domination \eqref{eq:slab-gaussian-domination}
with $h=h_{w'}$ and split the expectation into its sectors:
\[
 1\ge\sum_w
 \exp\bigl(\langle h_{w'},h_w\rangle-\|h_{w'}\|^2/2\bigr)P_n(w)
 \ge \exp(\|h_{w'}\|^2/2)P_n(w').
\]
Since $\|h_w\|^2=r_q^2|w|_2^2/n$, this proves the
first bound in \eqref{eq:transfer-flux-weight}.
The identity $\sum_wP_n(w)=1$ proves the second.

Set $x=\exp(-2\pi^2/(q^2n))$. Positivity permits the
factorization
\[
 \varsigma_n=\left(\sum_{i\in\Z}x^{i^2}\right)^3-1
 \le\left(1+\frac{2x}{1-x}\right)^3-1.
\]
For $q^2n\le1$, one has $x<2^{-18}$.
The bound $(1-x)^{-1}\le1+2x$ gives
$2x/(1-x)\le2.001x<10^{-5}$. Therefore
$\varsigma_n\le2.001x(3+3\cdot10^{-5}+10^{-10})
\le7x<1/8$. This also proves convergence of
\eqref{eq:transfer-flux-tail}.

Taking logarithms in \eqref{eq:transfer-zero-factor}
at $n$ and $jn$ gives \eqref{eq:transfer-defect-split}.
By \eqref{eq:external-photon-factor},
\[
 \log Z_n^{\mathrm{ph}}
   =(n^3+2)\log r_q+(n^3-1)\log(2\pi)-\mathcal L(n)/2.
\]
Its bulk terms cancel in the defect. The remaining
constant is $(j^3-1)\log(2\pi/q^2)$.
Equation \eqref{eq:transfer-photon-logdet} proves
\eqref{eq:transfer-photon-defect}. Finally
$\log P_{jn}(0)\le0$, and
$P_n(0)\ge1-\varsigma_n>0$ gives
\eqref{eq:transfer-flux-defect}.
\end{proof}

\section{Wick products on larger cutoff tori}\label{sec:vtm-wick}

Use the cell sources of Section~\ref{sec:reference-cells}
and write
\[
 N_*=8L^3,\qquad
 a_L=\frac{L^3-1}{(L-1)^3},\qquad
 D^{\mathrm{mat}}=L^3\log z_{r,n_*}-\log z_{r,Ln_*},
 \qquad
 \delta=\frac{D_L}{(L-1)^3N_*}.
\]
Fix $0<e\le1/3$, with $q=e\sqrt\eta$ and
$q^2n_*=e^2R^*$, and assume
\begin{equation}\label{eq:transfer-W3}
 6N_*e^2R^*\log(1/e)\le\pi^2.
\end{equation}
The choice of scales ensuring this condition is given
below. Let $\mathcal C_*=\{0,\ldots,2L-1\}^3$
and use the probe variance $v_\eta$ of
\eqref{eq:cells-orthogonality}. For real $c_2,c_4$ put
\begin{align}
 O_m^{[c]}&=W_m(X_c;v_\eta,c_m),\qquad m=2,4,
                                      \label{eq:transfer-Wick-observable}\\
 \Theta_m^{(0)}
  &=\left|
      \left[\prod_{c\in\mathcal C_*}
            (\partial_{t_c}^m-c_m)\right]\mathring\Phi(0)
    \right|^{1/N_*}.                  \label{eq:transfer-Theta}\\
 M_4&=600N_*^2v_\eta^2e^{-4}+|c_4|,\qquad
 M_2=12N_*v_\eta e^{-2}+|c_2|,\notag\\
 T_m&=\Theta_m^{(0)}+2e^6M_m,\qquad
 u_m=\exp(\delta)T_m.                 \label{eq:transfer-Wick-data}
\end{align}
Here $\mathring\Phi$ is the normalized external partition
function \eqref{eq:external-reference-partition}.
The notation $e^{-m}$ denotes an inverse power of the
charge.

\begin{proposition}\label{prop:transfer-Wick}
Assume the probe placement condition \eqref{eq:cells-fit}.
For $n=n_*$, or $n\in p_{\rm cell}\Z$ with $n\ge Ln_*$,
the all-flux law with the fixed real lattice mass obeys
\begin{align}
 0\le\delta
 &\le\frac{a_L}{N_*}\log\frac{2\pi}{e^2R^*}
       +\frac{K_L+D^{\mathrm{mat}}+1}{(L-1)^3N_*},
                                           \label{eq:transfer-delta}\\
 \left|\E_n^{\mathrm{af}}\prod_{c\in\mathcal C}
                        O_m^{[c]}\right|
 &\le u_m^{|\mathcal C|}              \label{eq:transfer-Wick-product}
\end{align}
for every set $\mathcal C$ of distinct cells and $m=2,4$.
In particular, if $\kappa_j$ denotes the cumulant of
$X_0$ on that torus, then
\begin{align}
 \left|e^{-2}(\kappa_2-v_\eta)-c_2\right|&\le u_2,
                                      \label{eq:transfer-second-cumulant}\\
 \left|e^{-4}\kappa_4-(c_4-3c_2^2)\right|
   &\le u_4+3u_2(2|c_2|+u_2).        \label{eq:transfer-fourth-cumulant}
\end{align}
The bounds are uniform in $n$.
Their right sides are formed on the two reference tori;
$\Theta_m^{(0)}$, and therefore $u_m$, may depend on
the placement.
\end{proposition}

\begin{proof}
Condition \eqref{eq:transfer-W3} implies $e^2R^*<1$,
since $N_*\ge1000$ and $\log(1/e)\ge\log3>1$.
Put $x=\exp(-2\pi^2/(e^2R^*))$. The same condition gives
\[
 x\le e^{12N_*},\qquad x^{1/(2N_*)}\le e^6.
\]
Apply Proposition~\ref{prop:transfer-flux-defect} at
$n_*$ and $j=L$. Since $\varsigma_{n_*}\le7x<1/8$,
\[
 D_L^{\mathrm{fl}}\le-L^3\log(1-\varsigma_{n_*})
              \le8L^3x=N_*x
              \le N_*3^{-12N_*}\le1.
\]
The photon estimate and the defect decomposition
now prove the upper bound in \eqref{eq:transfer-delta};
its lower bound is Proposition~\ref{prop:transfer-general}.

All the polynomials in \eqref{eq:transfer-Wick-observable}
are the same local integrands on every carrying torus,
by Section~\ref{sec:reference-cells}. They are real
elements of the cell observable class. Write
$O_m^\sharp=\prod_{c\in\mathcal C_*}O_m^{[c]}$ on
$C(n_*)$. Its all-flux expectation is nonnegative,
by the fold trace \eqref{eq:transfer-fold-trace}.
We estimate it by splitting flux zero from the
other sectors.

First, Gaussian domination and the moment bound
give, for every positive integer $i$,
\[
 \E_{n_*}^{\mathrm{af}}X_c^{2i}
                              \le(2i\exp(1)v_\eta)^i.
\]
Minkowski's inequality at exponent $2N_*$ yields
\begin{align*}
 \|H_4(X_c;v_\eta)\|_{2N_*}
   &\le[64\exp(2)N_*^2+24\exp(1)N_*+3]v_\eta^2
       \le600N_*^2v_\eta^2,\\
 \|H_2(X_c;v_\eta)\|_{2N_*}
   &\le[4\exp(1)N_*+1]v_\eta
       \le12N_*v_\eta.
\end{align*}
Thus $\|O_m^{[c]}\|_{2N_*}\le M_m$.
The finite moments justify all uses of H\"older
and Cauchy--Schwarz below, also when $v_\eta=0$.

On the zero-flux sector,
\eqref{eq:external-Wick-product} gives
\[
 \left|\E_{n_*}^{\mathrm{af}}
               [O_m^\sharp\mathbf1_{\{\mathrm{flux}=0\}}]\right|
   =P_{n_*}(0)\left|\E_{n_*}^{(0)}O_m^\sharp\right|
   \le(\Theta_m^{(0)})^{N_*}.
\]
On the remaining sectors, Cauchy--Schwarz followed
by H\"older with $N_*$ factors gives
\begin{align*}
 \left|\E_{n_*}^{\mathrm{af}}
               [O_m^\sharp\mathbf1_{\{\mathrm{flux}\ne0\}}]\right|
 &\le\varsigma_{n_*}^{1/2}
                \bigl(\E_{n_*}^{\mathrm{af}}|O_m^\sharp|^2\bigr)^{1/2}\\
 &\le\varsigma_{n_*}^{1/2}
                \prod_{c\in\mathcal C_*}\|O_m^{[c]}\|_{2N_*}\\
 &\le\bigl(\varsigma_{n_*}^{1/(2N_*)}M_m\bigr)^{N_*}.
\end{align*}
Since $\varsigma_{n_*}^{1/(2N_*)}\le2e^6$,
the nonnegative $N_*$-th root satisfies
\[
 \|O_m\|_*=
       \bigl(\E_{n_*}^{\mathrm{af}}O_m^\sharp\bigr)^{1/N_*}
             \le\Theta_m^{(0)}+2e^6M_m=T_m.
\]
For nonnegative $a,b$ we used
$(a^{N_*}+b^{N_*})^{1/N_*}\le a+b$.
Proposition~\ref{prop:transfer-general}, with
$s_0=2L$ and $j=L$, proves
\eqref{eq:transfer-Wick-product}.

For one cell the same estimate gives
$|e^{-m}\E H_m(X_0;v_\eta)-c_m|\le u_m$.
Charge conjugation makes the odd moments zero.
Writing $y=e^{-2}\E H_2(X_0;v_\eta)$, one has
\[
 e^{-4}\kappa_4=e^{-4}\E H_4(X_0;v_\eta)-3y^2.
\]
The second-order conclusion follows at once.
For the fourth, use $|y-c_2|\le u_2$ and
$|y+c_2|\le2|c_2|+u_2$ to obtain
\eqref{eq:transfer-fourth-cumulant}.
\end{proof}

\section{Reference scales and the loss of charge exponent}\label{sec:vtm-bounds}

The reference side is chosen from the logarithmic charge
scale. Use $L,M,h$ from the introduction and the localization
power two. Thus
\begin{equation}\label{eq:transfer-scale-choice}
 e_K=e\sqrt h,\qquad v_K=-\log e_K,\qquad
 r_K=\min\{L^i:L^i\ge v_K^2,\ i\in\Z_{\ge0}\},
 \qquad R^*=\frac{hMr_K}{L}.
\end{equation}
Assume $v_K\ge8$, $e\le1/3$, $e_K\log h\le1$, and
\begin{equation}\label{eq:transfer-tail-condition}
 6N_*Mv_K^2 \exp(-v_K)
                 \left(v_K\exp(-v_K)+\tfrac12\right)\le\pi^2.
\end{equation}
These are finite charge conditions after $h,L,M$ are fixed.
They hold on a sufficiently small charge interval:
$v_K\to\infty$ as $e\to0$, and the polynomial times
exponential in \eqref{eq:transfer-tail-condition}
tends to zero.

Minimality gives $v_K^2\le r_K<Lv_K^2$.
Consequently $T^*=\log_L(Mr_K)-1$ is an integer at
least one, $R^*=hL^{T^*}$, and
$\ell_c=R^*/(2L)\to\infty$ as $e\to0$.
Moreover $1/L<\bar mh\le1$, so
\begin{equation}\label{eq:transfer-volume-error}
 \bar mR^*>\frac{Mr_K}{L^2}\ge r_K\ge v_K,
 \qquad \exp(-\bar mR^*)\le e\sqrt h.
\end{equation}
The same upper bound on $r_K$ gives
\begin{align*}
 e^2R^*\log(1/e)
 &<Mv_K^2 \exp(-2v_K)
                          (v_K+\tfrac12\log h)\\
 &\le Mv_K^2 \exp(-v_K)
                          (v_K\exp(-v_K)+\tfrac12).
\end{align*}
Thus \eqref{eq:transfer-tail-condition} implies
\eqref{eq:transfer-W3}.
The probe fits every sufficiently small charge and
sufficiently fine mesh, since $D_\beta^{\mathrm{buf}}$
is fixed and $\ell_c\to\infty$.

We now isolate the two reference estimates used in the
passage to arbitrary volumes. Their dependence on the
charge matters: a bounded partition defect still loses
a small power of $e$ through its photon contribution.
Let $C_{\mathrm{mat}}$ and $C_{\mathrm{ref}}$ be finite constants
independent of the charge, mesh and placement. Assume
that for $N\ge N_{\mathrm{ref}}(e)$ the mass assignment and reference
response satisfy
\begin{align}
 D^{\mathrm{mat}}&\le C_{\mathrm{mat}},\notag\\
 \Theta_4^{(0)}&\le C_{\mathrm{ref}} e^{\mathfrak e},\qquad
 (1+|c_2|)\Theta_2^{(0)}
                          \le C_{\mathrm{ref}} e^{\mathfrak e},
                                   \label{eq:transfer-terminal-input}
\end{align}
where
\begin{equation}\label{eq:transfer-charge-exponents}
 \mathfrak e=\tfrac12(1/4-9\varepsilon)-12\varepsilon,\qquad
 \sigma=\mathfrak e-\frac{2a_L}{N_*}>0,\qquad
 0<\varepsilon\le1/200.
\end{equation}
Here $N_{\mathrm{ref}}(e)$ is a fixed finite mesh threshold.
These are the reference hypotheses of \cref{thm:vtm-main}.
For the common tuned mass, the matter estimate is
\cite[Theorem~A, (5)]{TerminalPressure}, and the response
estimate is \cite[Theorem~A]{ReferenceResponse}.
Their joint charge interval and cutoff threshold are part
of the construction data assumed here.

On the reference torus take the free coefficients
$c_m=N_0^{(m)}(0)$ from
\eqref{eq:response-coefficient-relations}, using the
source associated with the same placed probe.
\cite[Proposition~5.2]{PositiveProbe}
and the support fit give a constant, enlarged into
$C_{\mathrm{ref}}$ if necessary, such that
\begin{equation}\label{eq:transfer-free-input}
 |c_2|+|c_4|^{1/2}\le C_{\mathrm{ref}},\qquad
 |c_4-3c_2^2-b(f)|
              \le C_{\mathrm{ref}}(e+\eta^{1/4}).
\end{equation}
Indeed $R^*\ge D_\beta+1$, and
\eqref{eq:transfer-volume-error} bounds the finite-volume
error in \cite[Theorem~B]{PositiveProbe}.
The coefficients are real and independent of placement.

\begin{proposition}\label{prop:transfer-UV-inputs}
Fix the probe and the scale data above, and assume
\eqref{eq:cells-fit}, \eqref{eq:transfer-terminal-input} and \eqref{eq:transfer-free-input}.
There is a finite
constant $C_U$, independent of $e,N,T$ and placement,
such that for all the stated reference charges and
$N\ge N_{\mathrm{ref}}(e)$,
\begin{equation}\label{eq:transfer-small-Wick}
 u_4\le C_Ue^\sigma,\qquad
 (1+|c_2|)u_2\le C_Ue^\sigma,\qquad
 |c_4-3c_2^2-b(f)|\le C_U(e^\sigma+\eta^{1/4}).
\end{equation}
On every cutoff torus with $T\ge T^*+1$, every finite
set $Y\subset(R^*/L)\Z^3$ of points incongruent
modulo $R\Z^3$ satisfies
\begin{equation}\label{eq:transfer-translated-Wick}
 \left|\E_{N,T}\prod_{y\in Y}
       W_m(F^\sharp(\tau_yf^\eta);v_\eta,c_m)\right|
                          \le u_m^{|Y|},\qquad m=2,4,
\end{equation}
Here $W_m$ is the polynomial in \eqref{eq:vtm-Wick}.
The fourth cumulant of the fixed smooth test obeys
\begin{equation}\label{eq:transfer-smooth-cumulant}
 |\kappa_4(F^\sharp(f))-e^4b(f)|
 \le(8C_U+3C_U^2)e^{4+\sigma}
                  +C_U(\eta^2+e^4\eta^{1/4}).
\end{equation}
The tests $f^\eta$ have one fixed compact support,
$f^\eta\to f$ in $L^2$, and $v_\eta\to\|f\|_2^2$.
Thus \eqref{eq:transfer-small-Wick} and
\eqref{eq:transfer-translated-Wick} give exactly the
cutoff hypotheses of Proposition~\ref{prop:UV-component-cumulant}.
\end{proposition}

\begin{proof}
Write
\[
 s_\beta=\|f\|_2+\|\beta\|_{H^4},\qquad
 C_\Gamma^{\mathrm{tr}}=(2\pi)^{a_L/N_*}
         \exp\left(\frac{K_L+C_{\mathrm{mat}}+1}
                                  {(L-1)^3N_*}\right).
\]
Since $R^*\ge1$, \eqref{eq:transfer-delta} gives
\begin{equation}\label{eq:transfer-Gamma-loss}
 \exp(\delta)\le C_\Gamma^{\mathrm{tr}} e^{-2a_L/N_*}.
\end{equation}
The inverse power here is a power of the charge.
The source-curvature identity and Jensen's inequality
for the disjoint averaging cubes give
$v_\eta^{1/2}\le\|f^\eta\|_2\le s_\beta$.
Use this and \eqref{eq:transfer-free-input} in
\eqref{eq:transfer-Wick-data}:
\begin{align*}
 T_4
 &\le C_{\mathrm{ref}} e^{\mathfrak e}
           +1200N_*^2s_\beta^4e^2+2C_{\mathrm{ref}}^2e^6\\
 &\le(C_{\mathrm{ref}}+1200N_*^2s_\beta^4+2C_{\mathrm{ref}}^2)e^{\mathfrak e},\\
 (1+|c_2|)T_2
 &\le C_{\mathrm{ref}} e^{\mathfrak e}
      +2e^6(12N_*s_\beta^2e^{-2}+C_{\mathrm{ref}})(1+C_{\mathrm{ref}})\\
 &\le[C_{\mathrm{ref}}+2(12N_*s_\beta^2+C_{\mathrm{ref}})(1+C_{\mathrm{ref}})]
                                                    e^{\mathfrak e}.
\end{align*}
Here $0<\mathfrak e<1/4$, from
\eqref{eq:transfer-charge-exponents}.
One possible constant is
\begin{equation}\label{eq:transfer-CU}
 \begin{split}
 C_U=\max\{&
 C_\Gamma^{\mathrm{tr}}(C_{\mathrm{ref}}+1200N_*^2s_\beta^4+2C_{\mathrm{ref}}^2),\\
 &C_\Gamma^{\mathrm{tr}}[C_{\mathrm{ref}}+2(12N_*s_\beta^2+C_{\mathrm{ref}})(1+C_{\mathrm{ref}})],
              \ C_{\mathrm{ref}},\ 60s_\beta^4\}.
 \end{split}
\end{equation}
Multiplication by \eqref{eq:transfer-Gamma-loss}
proves the first two bounds in
\eqref{eq:transfer-small-Wick}. The last follows from
\eqref{eq:transfer-free-input}, since $e\le e^\sigma$.

For a cutoff side $n=8hL^{T+N}$ with $T\ge T^*+1$,
one has $n=8L^{T-T^*}n_*$, a multiple of $p_{\rm cell}$ with
$n\ge Ln_*$. For $y\in(R^*/L)\Z^3$ use the even
cell index $c(y)$ of
\eqref{eq:cells-translation-observable}.
Incongruent points give distinct cells.
Proposition~\ref{prop:transfer-Wick} applies to their
product, and lattice translation invariance removes
the common placement. This proves
\eqref{eq:transfer-translated-Wick}.
For a fixed finite set of distinct points it applies
as soon as $R>2\max_{y\in Y}|y|_\infty$.

It remains to replace the discrete smooth test
$f^\eta$ by $f$ in the cumulant. We give the
quantitative multilinear estimate used for this purpose.
For a real two-cochain $g$, charge conjugation and
Gaussian domination give
\[
 \E\cosh(t\langle g,F^{\mathrm u}\rangle)
                          \le\exp(t^2\|g\|^2/2).
\]
Extract its nonnegative term of degree $2i$ and choose
$t^2=2i/\|g\|^2$. This gives
\[
 \E\langle g,F^{\mathrm u}\rangle^{2i}
    \le(2i)!\left(\frac{\|g\|^2}{2i}\right)^i \exp(i).
\]
For $i=2$ the bound is $(3/2)\exp(2)\|g\|^4$.
H\"older bounds the four-field moment by
\[
 \tfrac32\exp(2)\prod_a\|g_a\|.
\]
The variance bound bounds each of its three pair subtractions
by $\prod_a\|g_a\|$. Since $(3/2)\exp(2)+3<15$,
the symmetric multilinear forms $K_2,K_4$ obey
\begin{equation}\label{eq:transfer-multilinear-cumulants}
 |K_2(g_1,g_2)|\le\|g_1\|\|g_2\|,\qquad
 |K_4(g_1,g_2,g_3,g_4)|\le15\prod_a\|g_a\|.
\end{equation}
For $g=0$ the corresponding moment is zero.

Let $g$ and $g'$ be the periodizations of the normalized
covering cochains $f_{[\eta]}$ and $(f^\eta)_{[\eta]}$
from \cref{sec:reference-cells} onto the cutoff torus.
The common support cube has side less than $R$.
The averaging bound and \eqref{eq:cells-test-error} give
\[
 \|g\|,\|g'\|\le s_\beta,\qquad
 \|g-g'\|\le\eta^2\|\beta\|_{H^4}.
\]
The observables paired with these cochains are
$F^\sharp(f)$ and $F^\sharp(f^\eta)$.
Telescope the four slots of $K_4$ and the two slots
of $K_2$ in \eqref{eq:transfer-multilinear-cumulants}.
This proves
\begin{align}
 |\kappa_4(F^\sharp(f))-\kappa_4(F^\sharp(f^\eta))|
                              &\le60s_\beta^4\eta^2,\notag\\
 |\kappa_2(F^\sharp(f))-\kappa_2(F^\sharp(f^\eta))|
                              &\le2s_\beta^2\eta^2.
                                   \label{eq:transfer-test-replacement}
\end{align}
One-cell transfer and \eqref{eq:transfer-small-Wick}
bound the fourth cumulant of $F^\sharp(f^\eta)$ by
\[
 |\kappa_4(F^\sharp(f^\eta))-e^4(c_4-3c_2^2)|
                     \le(7C_U+3C_U^2)e^{4+\sigma}.
\]
Adding the free-coefficient error and the first bound
in \eqref{eq:transfer-test-replacement} proves
\eqref{eq:transfer-smooth-cumulant}.
Finally the compact support, $L^2$ convergence and
variance convergence are respectively
\eqref{eq:cells-smooth-test}, \eqref{eq:cells-test-error}
and \eqref{eq:cells-variance-limit}.
\end{proof}


\section{Periodic fields and a prescribed continuum law}
\label{sec:continuum-limits}

A finite periodic step field does not satisfy the sharp global
$L^2(\R^3)$ Gaussian bound for arbitrary Schwartz tests.
The periodized norm must be retained until the volume tends
to infinity. We give the moment argument of
\cite[\S6.1--6.2]{QED3Master}; it is the step needed to use the
finite-volume product estimates for a prescribed parent.

Let $\mathsf A=1+|y|^2-\Delta$ on $\R^3$, with real
$L^2$-orthonormal Hermite two-forms $h_\kappa$ and eigenvalues
$\lambda_\kappa$. Each scalar Hermite function occurs once
for each of the three components. We use the Hilbert space
\[
 H_{-4}=\left\{F:\sum_\kappa\lambda_\kappa^{-8}
                                  |F(h_\kappa)|^2<\infty\right\}.
\]
It is continuously embedded in tempered two-form distributions,
and each Schwartz evaluation is continuous there.
The step-field realization, tightness and fixed-spacing limits
in this space are proved in \cite{ContinuumExistence}.

\subsection{The periodic seminorm}

For $R=n\eta$ and a Schwartz two-form $g$, set
$g^R(y)=\sum_{k\in\Z^3}g(y+Rk)$ and
\[
 g_\eta([\sigma])=\eta^{-3}
                \int_{Q_\eta(\sigma)}g^R_{\mu\nu}(y)\,\dd y,
 \qquad |g_\eta|_\eta^2=\eta^3\sum_{[\sigma]}|g_\eta([\sigma])|^2.
\]
Unfolding the absolutely convergent sum gives
\begin{equation}\label{eq:limit-step-pairing}
 F^\sharp(g)=\eta^{3/2}\sum_{[\sigma]}
                         g_\eta([\sigma])F^{\rm u}(\sigma).
\end{equation}
Cell Cauchy--Schwarz and Minkowski give
\begin{equation}\label{eq:limit-torus-norm}
 |g_\eta|_\eta\le\|g^R\|_{L^2(\R^3/R\Z^3)}\le p_R(g),
 \qquad
 p_R(g)=\|g\|_2+
      26R^{-2}\||y|^2g\|_{L^2(\R^3\setminus Q_R)},
\end{equation}
where $Q_R=[-R/2,R/2)^3$.
For completeness, put $a_k=\|g\|_{L^2(Q_R+Rk)}$.
The periodized norm is at most $\sum_ka_k$. On a cube with
$k\ne0$, $|y|\ge R|k|_\infty/2$, so
\[
 \sum_{k\ne0}a_k
 \le4R^{-2}\left(\sum_{k\ne0}|k|_\infty^{-4}\right)^{1/2}
                 \||y|^2g\|_{L^2(\R^3\setminus Q_R)}.
\]
The shell count $24j^2+2$ gives a sum
$4\pi^2+\pi^4/45<41.65$; its square root times four is
less than $26$. This proves \eqref{eq:limit-torus-norm}.
For $R\ge1$,
\[
 p_R(g)\le p_1(g),\qquad p_R(g)\longrightarrow\|g\|_2
                      \quad(R\longrightarrow\infty).
\]
The seminorm $p_1$ is continuous on Schwartz space.
If the support of $g$ lies in a translate of an open
$R$-cube, its distinct period translates are disjoint and
$|g_\eta|_\eta\le\|g\|_2$.

By \eqref{eq:slab-gaussian-domination},
\begin{equation}\label{eq:limit-cutoff-G}
 \E_{N,T}\exp(F^\sharp(g))
       \le\exp(|g_\eta|_\eta^2/2)\le\exp(p_R(g)^2/2)
                          \qquad(g\text{ real}).
\end{equation}

\begin{proposition}\label{prop:vtm-periodic-moments}
Suppose $N_i,T_i\to\infty$ and the laws $P_{N_i,T_i}$ of
$F^\sharp$ converge weakly on $H_{-4}$ to $\mu$.
Then $\E_\mu\exp(F(g))\le\exp(\|g\|_2^2/2)$ for every real
Schwartz two-form $g$. Every joint moment of a fixed finite
family of Schwartz tests converges along the same sequence.
\end{proposition}

\begin{proof}
Truncate $\exp(F(g))$ above, pass this bounded continuous
function to the limit, and remove the truncation.
Since $R_i\to\infty$, \eqref{eq:limit-torus-norm} proves the
Gaussian bound.

For fixed complex tests $g_1,\ldots,g_s$ put
\[
 Z_i(z)=\E_{N_i,T_i}\exp\left(\sum_a z_aF^\sharp(g_a)\right),
 \qquad
 Z(z)=\E_\mu\exp\left(\sum_a z_aF(g_a)\right).
\]
The evaluations are continuous on $H_{-4}$, and
\[
 \E_{N_i,T_i}
 \left|\exp\left(\sum_a z_aF^\sharp(g_a)\right)\right|^2
 \le\exp\left(2p_1\left(\operatorname{Re}\sum_a z_ag_a\right)^2\right).
\]
This is uniform in $i$ and on compact $z$-sets.
Uniform integrability and weak convergence give pointwise
convergence of $Z_i$ to $Z$. The exponential bounds make
these functions entire and locally uniformly bounded.
Cauchy's estimates give equicontinuity on compact sets;
pointwise convergence is therefore locally uniform.
Cauchy's formula passes every derivative at zero to the limit,
which proves joint moment convergence.
\end{proof}

We also record the varying-test estimate used below.
If $g_1,\ldots,g_J$ are real and their supports lie in one
open period cube, the Gaussian moment bounds and H\"older
give
\[
 \E\prod_{a=1}^J|F^\sharp(g_a)|
                     \le c_J\prod_{a=1}^J\|g_a\|_2
\]
for a constant $c_J$ depending only on $J$.
The same estimate holds for $\mu$.
Thus replacing each factor from a bounded family of tests
by another test, with $L^2$ difference at most $\delta$,
changes the expectation of a degree-$J$ monomial by at most
$Jc_JB^{J-1}\delta$, where $B$ bounds their $L^2$ norms.
This follows by telescoping the $J$ slots. For any fixed
finite family of compactly supported translates, the
support condition holds once $R_i$ is sufficiently large.

\subsection{Realizing the chosen parent by finite tori}

Let $\mu_i$ and $\mu$ be the prescribed laws in
\cref{thm:vtm-main}. Fix a metric $d_{\rm w}$ for weak
convergence of probability measures on the separable Hilbert
space $H_{-4}$. At spacing $\eta_i=L^{-N_i}$, take the
defining volume sequence for $\mu_i$. Choose an element
$T_i$ of that sequence so large that
\[
 d_{\rm w}(P_{N_i,T_i},\mu_i)<1/i,\qquad
 T_i\ge i+T^*(e)+1.
\]
Then
\[
 d_{\rm w}(P_{N_i,T_i},\mu)
 \le1/i+d_{\rm w}(\mu_i,\mu)\longrightarrow0.
\]
Both cutoffs diverge. This is the prescribed-parent
diagonal of \cite[Corollary~7.4]{ContinuumExistence},
with the finite-periodic moment conclusion supplied by
\cref{prop:vtm-periodic-moments}. Discard finitely many terms
so that $N_i\ge N_{\rm ref}(e)$.
At this fixed charge the translation spacing
$\lambda=R^*(e)/L$ is independent of $i$.

The four reference quantities
$c_{2,i},c_{4,i},u_{2,i},u_{4,i}$ are bounded by the estimates
of \cref{prop:transfer-UV-inputs}. A subsequence on which
they converge still has weak limit $\mu$.
For every fixed finite $Y\subset\lambda\Z^3$, the points
are eventually incongruent modulo the torus and all the
translated supports fit in a single open period cube.
The next section passes the Wick products along this
same coefficient subsequence. The threshold for support
containment may depend on $Y$; the coefficients and parent
law do not.

\section{Fourth cumulants of the limiting components}
\label{sec:ultraviolet-limits}

The fourth-cumulant estimate for a mixed limiting law does not by
itself control each ergodic component. We use the stronger estimate
on products of translated Wick polynomials. It bounds their
translation-invariant conditional expectations, and hence the
second and fourth moments of almost every component.

\subsection{From distinct-point products to the invariant part}

Let $U(y)$,
$y\in\Z^d$, be the operators of a measure-preserving action on
a probability space, where $d\ge1$, with $U(y)G=G\circ T_{-y}$. Let
$P_{\mathrm{lat}}$ be the orthogonal projection in $L^2$ onto
the vectors invariant under this action.

\begin{lemma}\label{lem:UV-invariant-product-bound}
Suppose that $G$ is real and belongs to every $L^p$, $p<\infty$.
If $u\ge0$ and
\begin{equation}\label{eq:UV-distinct-product-bound}
 \left|\E\prod_{a=1}^J U(y_a)G\right|\le u^J
\end{equation}
for every finite family of pairwise distinct points of $\Z^d$,
then $|P_{\mathrm{lat}}G|\le u$ almost surely.
\end{lemma}

\begin{proof}
For $B_n=\{0,\ldots,n-1\}^d$ put
$A_n=n^{-d}\sum_{y\in B_n}U(y)$. These contractions converge
strongly in $L^2$ to $P_{\mathrm{lat}}$.
Indeed, the orthogonal complement of the invariant space is the
closed span of $(U(y)-1)H$. On such a vector, cancellation of
the common terms in $B_n$ and $y+B_n$ gives
\[
 \|A_n(U(y)-1)H\|_2\le2|y|_1n^{-1}\|H\|_2.
\]
On invariant vectors the averages are the identity, proving
the asserted convergence by the orthogonal decomposition.

For an even integer $J\ge2$, expand $\E(A_nG)^J$.
The distinct-index terms obey \eqref{eq:UV-distinct-product-bound}.
There are at most $\binom J2 n^{d(J-1)}$ remaining tuples:
each has a repeated pair of indices. H\"older's inequality
bounds their individual expectations in modulus by
$\|G\|_J^J$. Therefore
\[
 \E|A_nG|^J
       \le u^J+\binom J2 n^{-d}\|G\|_J^J.
\]
Along a subsequence the $L^2$ convergence is almost sure.
Fatou's lemma gives $\E|P_{\mathrm{lat}}G|^J\le u^J$
for every even $J$. For $t>u$, Markov's inequality now gives
$\Pr(|P_{\mathrm{lat}}G|>t)\le(u/t)^J$.
Let $J\to\infty$ and then $t\downarrow u$ through a countable
sequence.
\end{proof}

Let a translation-invariant law $\mu$ on tempered fields have
the decomposition $\mu=\int\mu_\xi\,\dd\nu(\xi)$.
Write $I$ for its translation-invariant sigma algebra and
$P=\E_\mu[\,\cdot\mid I]$. Apply the lemma to translations
on $\lambda\Z^3$, with $\lambda>0$ fixed. The range of $P$
is contained in that of $P_{\mathrm{lat}}$, so
$P=PP_{\mathrm{lat}}$. Positivity of conditional expectation
therefore gives
\begin{equation}\label{eq:UV-component-product-bound}
 |PG|\le P|P_{\mathrm{lat}}G|\le u,\qquad
 |\mu_\xi(G)|\le u\quad\text{for }\nu\text{-almost every }\xi.
\end{equation}
The second assertion uses the conditional-law representation of
the decomposition. Integrability under the conditional laws and
the representation for this integrable $G$ follow by truncating
its positive and negative parts.

\subsection{Passing the Wick-product estimate to a limit law}

Fix the charge $e>0$ and a real compactly supported test $f$.
Use the Wick polynomials $W_m$ in \eqref{eq:vtm-Wick},
with variance argument $v\ge0$ and $m=2,4$.

Take a joint cutoff sequence defining a limit law $\mu$.
Suppose that the lattice estimates give compactly
supported real tests $f_i$, all supported in one fixed cube,
and real numbers
\[
 v_i,\ c_{2,i},\ c_{4,i},\ u_{2,i},\ u_{4,i},
\]
with
$v_i\ge0$, such that
\begin{equation}\label{eq:UV-test-convergence}
 \|f_i-f\|_2\to0,\qquad
 v_i\to\|f\|_2^2,\qquad
 \sup_i(|c_{2,i}|+|c_{4,i}|)<\infty.
\end{equation}
For a fixed $\lambda>0$, assume that
\begin{equation}\label{eq:UV-cutoff-product-bound}
 \left|\E_i\prod_{y\in Y}
       W_m(F^\sharp(\tau_y f_i);v_i,c_{m,i})\right|
                     \le u_{m,i}^{|Y|},\qquad m=2,4,
\end{equation}
for every finite $Y\subset\lambda\Z^3$ and all sufficiently
large $i$ depending on $Y$. This is the translated-cell form
of the transfer estimate; the volume condition makes the
points of $Y$ incongruent modulo the torus.
The reference estimates also give
\begin{equation}\label{eq:UV-reference-bounds}
 \begin{split}
 0\le u_{4,i}&\le Ce^\sigma,\qquad
 0\le(1+|c_{2,i}|)u_{2,i}\le Ce^\sigma,\\
 |c_{4,i}-3c_{2,i}^2-b(f)|
                  &\le C(e^\sigma+\eta_i^{1/4}),
 \end{split}
\end{equation}
with $C$ independent of the cutoff, $\sigma>0$ and $0<e\le1$.
For the chamber laws, Proposition~\ref{prop:transfer-UV-inputs}
supplies these hypotheses throughout the charge interval
\eqref{eq:constants-probe-charge}.

Pass to a subsequence on which the four bounded sequences
$c_{2,i},c_{4,i},u_{2,i},u_{4,i}$ converge, with limits
$c_2,c_4,u_2,u_4$. The weak limit remains $\mu$.
Equation \eqref{eq:UV-reference-bounds} gives
\begin{equation}\label{eq:UV-reference-limit}
 u_4\le Ce^\sigma,\quad
 (1+|c_2|)u_2\le Ce^\sigma,\quad
 |c_4-3c_2^2-b(f)|\le Ce^\sigma.
\end{equation}
Set $G_m=W_m(F(f);\|f\|_2^2,c_m)$.
Then \eqref{eq:UV-cutoff-product-bound} passes to
\begin{equation}\label{eq:UV-law-product-bound}
 \left|\E_\mu\prod_{y\in Y}U(y)G_m\right|\le u_m^{|Y|}
                   \qquad(Y\subset\lambda\Z^3\text{ finite}).
\end{equation}

To justify this passage for the varying tests, expand the finite
product of Wick polynomials into monomials. For a monomial
of positive total degree $J$, telescope the replacement of each
$F^\sharp(\tau_y f_i)$ by $F^\sharp(\tau_y f)$.
All these tests have support in one open $R_i$-cube once $i$
is large. The test-map variance and product-moment bounds
then give an error at most
\[
 Jc_J B^{J-1}\|f_i-f\|_2,
 \qquad B=\sup_i\max(\|f_i\|_2,\|f\|_2).
\]
This tends to zero. The moments of the fixed tests converge
by Section~\ref{sec:continuum-limits}. The polynomial coefficients
converge by the selected coefficient subsequence and the
variance convergence in \eqref{eq:UV-test-convergence}.
There are only finitely many monomials for each fixed $Y$.
This proves \eqref{eq:UV-law-product-bound}.

\subsection{The cumulant in an ergodic component}

\begin{proposition}\label{prop:UV-component-cumulant}
Under \eqref{eq:UV-test-convergence}--\eqref{eq:UV-reference-bounds},
suppose that almost every translation-ergodic component of $\mu$
satisfies Gaussian domination. Then
\begin{equation}\label{eq:UV-component-cumulant}
 \left|\kappa_4^{\mu_\xi}(f)-e^4b(f)\right|
                 \le(8C+3C^2)e^{4+\sigma}
\end{equation}
for almost every component. In particular, if $b(f)>0$ and
$(8C+3C^2)e^\sigma\le b(f)/2$, its fourth cumulant is positive.
\end{proposition}

\begin{proof}
Gaussian domination of $\mu$ makes the real polynomials $G_2,G_4$
belong to every $L^p(\mu)$. Apply
Lemma~\ref{lem:UV-invariant-product-bound} to
\eqref{eq:UV-law-product-bound}, and then use
\eqref{eq:UV-component-product-bound}. On one set of full
component measure,
\[
 |d_m|\le u_m,\qquad d_m=\E_{\mu_\xi}G_m,\qquad m=2,4.
\]
Intersect with the full-measure set of components having Gaussian
domination. For a component in this intersection, put $X=F(f)$
and $v=\|f\|_2^2$. Gaussian domination gives $\E_{\mu_\xi}X=0$
and all moments finite. Thus
\[
 \kappa_4^{\mu_\xi}(f)
  =\E_{\mu_\xi}H_4(X;v)
          -3\bigl(\E_{\mu_\xi}H_2(X;v)\bigr)^2.
\]
The definition of $G_m$ gives the exact identity
\begin{equation}\label{eq:UV-cumulant-identity}
 e^{-4}\kappa_4^{\mu_\xi}(f)
                =c_4+d_4-3(c_2+d_2)^2.
\end{equation}
Consequently \eqref{eq:UV-reference-limit} implies
\begin{align*}
 |e^{-4}\kappa_4^{\mu_\xi}(f)-b(f)|
 &\le |c_4-3c_2^2-b(f)|+u_4+6|c_2|u_2+3u_2^2\\
 &\le8Ce^\sigma+3C^2e^{2\sigma}
 \le(8C+3C^2)e^\sigma.
\end{align*}
This proves the assertion uniformly over the full-measure set
of components. A centered normal variable has zero fourth
cumulant, so positivity also proves that these laws are
non-Gaussian.
\end{proof}




\subsection{Completion for the prescribed volume-first law}

\begin{proof}[Proof of \Cref{thm:vtm-main}]
The scale calculation in \cref{sec:vtm-bounds} gives
\eqref{eq:transfer-W3} and the reference volume error.
The reference and free-coefficient hypotheses are exactly
\eqref{eq:transfer-terminal-input} and \eqref{eq:transfer-free-input}.
Thus \cref{prop:transfer-UV-inputs}, with the uniform constant
\eqref{eq:transfer-CU}, proves
\eqref{eq:vtm-main-errors}--\eqref{eq:vtm-main-finite-cumulant}.

For the prescribed parent, use the finite-torus diagonal
in \cref{sec:continuum-limits}. Set
$f_i=f^{\eta_i}$ and $v_i=v_{\eta_i}$.
The compact support, test convergence and variance convergence
are \eqref{eq:cells-smooth-test},
\eqref{eq:cells-test-error} and \eqref{eq:cells-variance-limit}.
At this fixed charge, $\lambda=R^*/L$ is fixed.
Each finite $Y\subset\lambda\Z^3$ eventually consists of
distinct torus translates. The four reference sequences have
a convergent subsequence, whose weak limit is still the
prescribed $\mu$. The product passage proved above therefore
applies to this same law, with $C=C_U$.
Its components have Gaussian domination by the explicit
hypothesis \eqref{eq:vtm-component-G}.
Apply \cref{prop:UV-component-cumulant} and impose
\eqref{eq:constants-probe-charge} when $b(f)>0$.

If a separate argument gives other properties on a full set
of components of this parent, intersect that set with the
one obtained here. The component laws themselves have not
been replaced.
\end{proof}
\printbibliography
\end{document}
